% !TEX TS-program = xelatex
% !TEX encoding = UTF-8 Unicode
% !Mode:: "TeX:UTF-8"

\documentclass{resume}
\usepackage{zh_CN-Adobefonts_external} % Simplified Chinese Support using external fonts (./fonts/zh_CN-Adobe/)
%\usepackage{zh_CN-Adobefonts_internal} % Simplified Chinese Support using system fonts
\usepackage{linespacing_fix} % disable extra space before next section
\usepackage{cite}
\usepackage{graphicx}
% \usepackage{multirow}
% \usepackage{tabu}

\begin{document}
\pagenumbering{gobble} % suppress displaying page number

\name{张钏楠}

\basicInfo{
  \email{edwardzcn98@gmail.com} \textperiodcentered\ 
  \phone{(+86) 15630152367} \textperiodcentered\ 
  \linkedin[Edwardzcn]{https://www.linkedin.com/in/edwardzcn}
%   \github[Edwardzcn]{https://github.com/Edwardzcn} \textperiodcentered\
  % \instagram[Edwardzcn]{https://www.instagram.com/edwardzcn/} \textperiodcentered
  % \homepage[edwardzcn98yx.com]{https://www.edwardzcn98yx.com/}
}
 
\section{\faGraduationCap\  教育背景}
\datedsubsection {\textbf{中国科学技术大学}}{ 2021年9月 - 2026年6月}
\textit{硕士}\ 计算机系统结构 \ 计算机学院
\datedsubsection {\textbf{中南大学}} {2017年9月 - 2021年6月}
\textit{学士}\ 计算机科学与技术 \ 计算机学院 


\section{\faBriefcase\  工作经历}

% \datedsubsection{\textbf{中国科学技术大学先进数据系统实验室}\ 合肥, 安徽}{2021年2月 - 2021年6月}
% %\role{}
% % \begin{onehalfspacing}
% {(分布式一致性) 研究助理实习}
% \begin{itemize}
%   \item 设计实现并验证基于EPaxos的分布式一致性协调系统
%   \item 学习关于分布式系统、形式化验证相关知识
% \end{itemize}
% % \end{onehalfspacing}

\datedsubsection{\textbf{新加坡国立大学 TEST Lab}\ 新加坡}{2023年8月 - 2023年12月}
{（数据库测试）学术访问实习}
\begin{itemize}
  \item 为开源数据库测试框架 SQLancer 增加 SurrealDB 的支持，包括语法生成器以及对应PQS、TLP测试预言模块。
  \item 基于 SQLancer 开发针对跨数据库数据变更捕获工具的自动化测试套件。
\end{itemize}


\datedsubsection{\textbf{阿里巴巴PolarDB团队}\ 杭州}{2022年8月 - 2023年7月}
{（分布式数据库）研发实习}
\begin{itemize}
  \item 协助构建一个使用 Multi-Paxos 协议的地理分布式多节点读写数据库，实现行级多写。
  \item 基于 Sysbench 基准测试套件开发自动化生成跨地域节点性能测试工具。
\end{itemize}

\datedsubsection{\textbf{华为上海研究院OS内核实验室}\ 上海}{2020年10月 - 2021年1月}
%\role{Golang, Linux}
% \begin{onehalfspacing}
{（形式化验证）研发实习}
\begin{itemize}
  \item 使用ocamllex、ocamlyacc工具协助团队构建Lustre V6编译器项目前端词法语法部分
\end{itemize}
% \end{onehalfspacing}


\section{\faCogs 竞赛经历}


\datedsubsection{\textbf{2020年全国高校密码数学挑战赛}\ 北京}{2020年8月}
% \begin{onehalfspacing}
{二等奖（全国第三名）}
\begin{itemize}
  \item 针对低轮 Keccak256 哈希算法原像攻击，负责算法研究和编程实现
  \item 复现 Asiacrypt 会议的攻击方法，利用 NTL 数学库加速求解过程，改良特化结构进行原像攻击
\end{itemize}
% \end{onehalfspacing}

\datedsubsection{\textbf{2019年中国大学生程序竞赛-全国邀请赛（湖南）}\ 湘潭}{2019年5月}
{铜奖}
\begin{itemize}
  \item 参加湘潭大学主办的 CCPC 线下赛，获得铜奖
\end{itemize}

% Reference Test
%\datedsubsection{\textbf{Paper Title\cite{zaharia2012resilient}}}{May. 2015}
%An xxx optimized for xxx\cite{verma2015large}
%\begin{itemize}
%  \item main contribution
%\end{itemize}

\section{\faBook\ 教学经历}
\begin{itemize}
  \item \textbf{离散数学}课程助教 \ 2020年秋季学期 \ 中南大学 
  \item \textbf{编译原理}课程助教 \ 2021年春季学期 \ 中南大学
  \item \textbf{编译原理}课程助教 \ 2022年春季学期 \ 中国科学技术大学
  \item \textbf{全国大学生计算机系统能力大赛编译系统设计赛}指导老师 \ 2022年 \ 中国科学技术大学
\end{itemize}

\section{\faUsers\ 社团和组织经历}
\begin{itemize}
    \item 曾任中南大学信息学院团学会宣传部部长，负责歌手大赛、新生入学等线上线下宣传工作
    \item 中南大学ACM-ICPC算法竞赛实验室队员，协助管理实验室事务
    % \item 中南大学disc0ver开源组织主要贡献者，搭建知识Wiki，担任两门专业核心课助教
    
\end{itemize}
\section{\faTasks\ 技能与爱好}
% increase linespacing [parsep=0.5ex]
\begin{itemize}[parsep=0.4ex]
  \item \pbar{语言}{0.5}: 英语-流畅， 汉语-母语
  \item \pbar{文体}{0.8}: 摄影，电子琴，吉他，长跑，徒步登山（校登山队成员）
\end{itemize}

% \section{\faHeartO\ 获奖情况}
% \datedline{\textit{第一名}, xxx 比赛}{2013 年6 月}
% \datedline{其他奖项}{2015}

% \section{\faInfo\ 其他}
% % increase linespacing [parsep=0.5ex]
% \begin{itemize}[parsep=0.5ex]
%   \item 技术博客: http://blog.yours.me 
%   \item GitHub: https://github.com/username
%   \item 语言: 英语 - 熟练(TOEFL xxx)
% \end{itemize}

% %% Reference
% %\newpage
% %\bibliographystyle{IEEETran}
% %\bibliography{mycite}
\end{document}
